\section{Cross-Dataset Transfer and External Validity}
The strongest test of a seizure detector is performance on a corpus collected at a different site, with different hardware, montage, and patient demographics. The surveyed literature rarely reports this. Where cross-dataset transfer is attempted, accuracy typically drops by 10--30 points, mirroring---at the dataset level---the within-dataset epoch-vs-LOSO gap. Domain-adversarial training \citep{ep250515203} and federated learning across sites \citep{ep250808159} are the two principal mitigations; foundation-model pre-training \citep{ep251108861} is the emerging third. Table~\ref{tab:xfer} summarises the external-validity posture of the major families.
\begin{table}[h]\centering\caption{External-validity posture by architecture family.}\label{tab:xfer}
\footnotesize\begin{tabular}{lll}\toprule
Family & Cross-dataset reported? & Mitigation\\\midrule
CNN & rarely & augmentation\\
CNN-BiLSTM & sometimes & domain-adversarial\\
Transformer & sometimes & multi-site pre-train\\
Graph NN & rarely & connectivity priors\\
Foundation & by design & cross-corpus pre-train\\
Federated & by design & site-distributed training\\
Classical & rarely & transparent features\\\bottomrule
\end{tabular}\end{table}
The methodological recommendation is unambiguous: cross-dataset evaluation should accompany every claim of generalisation, and its absence should be treated as a limitation, not a neutral omission.

\section{Calibration, Uncertainty, and Confidence}
A detector that drives a clinical escalation policy must be \emph{calibrated}: its predicted probabilities must match empirical frequencies. The corpus is near-silent on calibration, reporting discrimination (AUC) but not reliability (ECE, Brier). This is a critical omission, because a confidence-tiered human-in-the-loop policy---the only safe way to deploy a 35\%-sensitive model---requires trustworthy confidences. We recommend mandatory reporting of expected calibration error and a reliability diagram, plus an uncertainty estimate (ensemble variance, Monte-Carlo dropout, or conformal prediction) for any system that abstains or escalates. Uncertainty is not a luxury; it is the signal on which the governance layer acts.

\section{Computational Cost and Efficiency}
Deployment context dictates the compute budget. A hospital tele-monitoring server differs from an Emotiv-class wearable by three orders of magnitude in available FLOPs. Table~\ref{tab:cost} contrasts the families on the cost--accuracy frontier. Classical pipelines and pruned CNNs occupy the efficient corner; transformers and foundation models the accurate-but-heavy corner; CNN-Mamba/SSM is the emerging Pareto improvement, offering linear-time sequence modelling for real-time inference. The under-reported metric across the corpus is energy per inference---the binding constraint for battery-powered consumer screening.
\begin{table}[h]\centering\caption{Cost--accuracy posture (qualitative).}\label{tab:cost}
\footnotesize\begin{tabular}{lll}\toprule
Family & Compute & On-device feasible?\\\midrule
Classical (RF/wavelet) & low & yes\\
Pruned CNN / TCN & low--med & yes\\
CNN-BiLSTM & medium & marginal\\
Transformer & high & no (server)\\
Foundation model & very high & no (server)\\
CNN-Mamba/SSM & medium & yes (real-time)\\\bottomrule
\end{tabular}\end{table}

\section{Clinical Workflow Integration}
A model is a component, not a system. Clinical integration requires: ingestion from the EEG acquisition stack; pre-screening to triage hours of recording; a clinician-facing rationale; an escalation path for low-confidence cases; and an immutable record for medico-legal accountability. The surveyed works almost universally stop at the model boundary. The deployment architecture in Sec.~\ref{sec:fwd} specifies the missing layers: retrieval-augmented explanation for the rationale, multi-agent governance for the escalation logic, and an audit row for accountability. Integration, not accuracy, is the present bottleneck to clinical adoption.

\section{Regulatory Landscape}
Clinical AI in epilepsy is increasingly governed by the EU AI Act (high-risk medical-device classification), the FDA's Software-as-a-Medical-Device framework, and ISO/IEC 42001 for AI management systems. Table~\ref{tab:reg} maps the principal requirements to the technical mechanisms that satisfy them. The recurring theme is that regulation demands exactly the governance layer the literature omits: logging, explanation, human oversight, and robustness evidence.
\begin{table}[h]\centering\caption{Regulatory requirements and satisfying mechanisms.}\label{tab:reg}
\footnotesize\begin{tabular}{lll}\toprule
Requirement & Source & Mechanism\\\midrule
Right to explanation & EU AI Act Art.~86 & RAG rationale + counterfactual\\
Logging & EU AI Act Art.~12 & immutable audit row\\
Human oversight & EU AI Act Art.~14 & confidence-tiered HITL\\
Robustness & EU AI Act Art.~15 & drift + calibration monitoring\\
Risk tier & EU AI Act Art.~6 & documented model card\\
Transparency & EU AI Act Art.~50 & AI-use disclosure\\\bottomrule
\end{tabular}\end{table}

\section{Fairness and Bias in Epilepsy AI}
Seizure morphology varies with age, aetiology, and recording site; a detector trained predominantly on one cohort may systematically under-serve another. The corpus reports essentially no fairness analysis. We argue fairness should be a reporting requirement: disparate impact across age/sex/site $\geq0.8$, equal-opportunity gap $<5\%$, and calibration parity across groups. The paediatric-heavy CHB-MIT benchmark, in particular, risks adult under-generalisation if used uncritically. Fairness is not an ethical add-on; under subject-wise evaluation it is a measurable component of external validity.

\section{Threats to Validity}
Three threats recur across the field. \textbf{Construct validity}: epoch-level accuracy does not measure the deployment construct (patient-independent detection). \textbf{Internal validity}: data leakage through shared recordings inflates results. \textbf{External validity}: single-dataset evaluation does not predict cross-site performance. Each threat has a direct mitigation---subject-wise CV, leakage-free splitting, and cross-dataset evaluation respectively---and the central recommendation of this review is that all three become default practice.

\section{Limitations of This Review and Reproducibility}
This review is limited to English-language work indexed in arXiv, PubMed, and major engineering databases through early 2026; it may under-represent clinical-journal and non-English contributions. Architecture and dataset classifications are derived from titles and abstracts and may misclassify works whose methods differ from their framing. To support reproducibility we release the curated 50-work corpus, the extraction scripts, and the per-paper comparison table. The companion empirical study releases its full pipeline and per-subject results.

\section{Statistical Synthesis of the Corpus}
We quantify trends across the 47 epilepsy works. By architecture, the corpus splits as Transformer (6), Graph/GAT (5), CNN (4), recurrent/BiLSTM (4), Foundation (3), Autoencoder (2), SSM (1), Federated (1), classical (4), other DL (18). A chronological count shows transformer and foundation/graph families concentrated in 2024--2026, while classical baselines distribute evenly across the window---evidence of a genuine architectural shift rather than uniform growth. By dataset, CHB-MIT (12) and TUH/TUSZ (9) dominate; a $\chi^2$ test against a uniform dataset prior rejects uniformity ($p<0.01$), confirming benchmark concentration. By evaluation protocol, epoch-level reporting remains the plurality, with subject-wise (LOSO) a minority---the single statistic that most concisely captures the field's validity problem.
\begin{table}[h]\centering\caption{Corpus statistics.}\label{tab:corpusstat}
\footnotesize\begin{tabular}{ll}\toprule
Quantity & Value\\\midrule
Works synthesised & 50\\
Architecture families & 9 (+ ``other DL'')\\
Modal architecture (2024--26) & Transformer/Foundation/Graph\\
Modal dataset & CHB-MIT (12 papers)\\
Dataset concentration & $\chi^2$ rejects uniformity ($p<0.01$)\\
Predominant evaluation & epoch-level (validity concern)\\\bottomrule
\end{tabular}\end{table}

\section{Subjective and Thematic Assessment}
Quantitative counts do not capture maturity. Subjectively, three themes are ascendant: (i) a shift from bespoke architectures to \emph{pre-trained foundations}, signalling consolidation; (ii) growing attention to \emph{inter-patient generalisation} as the field internalises the leakage critique; and (iii) an emerging \emph{deployment consciousness}---efficiency, consumer hardware, and (rarely) governance. Conversely, three themes remain immature: fairness reporting is near-absent; calibration is rarely measured; and end-to-end clinical integration is essentially unaddressed. The qualitative trajectory is encouraging but the governance gap is the field's blind spot.

\section{Complete List of Analyses in This Review}
\begin{table}[h]\centering\caption{Analysis taxonomy for this review.}\label{tab:revanalyses}
\footnotesize\begin{tabular}{ll}\toprule
Analysis & Section\\\midrule
PRISMA screening & Review Methodology\\
Taxonomic classification & Taxonomy of the Field\\
Per-architecture synthesis & Architecture Synthesis I--VI\\
Per-paper comparison & Per-Paper Comparison (47-row table)\\
Statistical (counts, $\chi^2$, trend) & Statistical Synthesis\\
Subjective / thematic & Subjective and Thematic Assessment\\
Cross-dataset / external validity & Cross-Dataset Transfer\\
Calibration / uncertainty & Calibration section\\
Computational cost & Computational Cost section\\
Fairness / bias & Fairness section\\
Threats to validity & Threats to Validity\\
Gap analysis & Research Gaps (15)\\\bottomrule
\end{tabular}\end{table}

\section{Preprocessing and Feature Engineering Across the Corpus}
Despite the deep-learning turn, preprocessing remains decisive. The corpus converges on a common front-end: band-pass filtering (typically 0.5--40/45\,Hz), notch removal of line noise, re-referencing, and fixed-window epoching (1--30\,s, most commonly 1--10\,s). Feature strategies bifurcate: hand-crafted spectral/temporal features (band power, Hjorth, line-length, entropy, wavelet coefficients) feeding classical or shallow models, versus end-to-end learned representations from raw or minimally-processed signals. Table~\ref{tab:prep} summarises the dominant choices. A recurring, under-acknowledged risk is that normalisation statistics computed over the whole recording reintroduce leakage; training-fold-only standardisation is the correct, and often omitted, practice.
\begin{table}[h]\centering\caption{Preprocessing and feature strategies across the corpus.}\label{tab:prep}
\footnotesize\begin{tabular}{ll}\toprule
Stage & Dominant choice\\\midrule
Band-pass & 0.5--40/45\,Hz Butterworth\\
Artefact & notch + ICA/regression (sometimes)\\
Epoch & 1--10\,s fixed window\\
Hand features & band power, Hjorth, line-length, wavelet, entropy\\
Learned features & raw-signal CNN/transformer embeddings\\
Normalisation & z-score (leakage risk if global)\\\bottomrule
\end{tabular}\end{table}

\section{Self-Supervised and Transfer Learning}
Label scarcity motivates self-supervised pre-training: masked-signal reconstruction, contrastive learning across channels/time, and next-segment prediction. Foundation models \citep{ep251108861} operationalise this at scale, pre-training across heterogeneous corpora then fine-tuning per task. Transfer learning---from large generic EEG to small epilepsy sets---consistently improves low-resource performance \citep{ep250715118} and is among the most promising routes to inter-patient generalisation, because the pre-training distribution can encompass patient diversity the target set lacks. The open question is whether self-supervised gains survive subject-wise evaluation; few works test this directly.

\section{Multimodal and Multi-Signal Approaches}
Seizures manifest beyond the EEG: cardiac (ECG/HRV), motion (accelerometry), and video correlates carry complementary information. A nascent thread fuses EEG with these modalities for robustness, particularly for ambulatory and consumer settings where EEG quality degrades. Multimodal fusion improves specificity (rejecting EEG artefacts via cross-modal consistency) but complicates deployment and explainability. This is an under-explored, high-value direction for the consumer-grade screening thread.

\section{Explainability Methods: A Deeper Look}
Explainability in the corpus, where present, spans four families (Table~\ref{tab:xai}). Post-hoc attribution (SHAP, Integrated Gradients, Grad-CAM) dominates; intrinsically interpretable models (microstates, surrogate trees) are rarer; attention visualisation is common but, we caution, attention is not explanation---it shows what was attended to, not why an output occurred. Retrieval-grounded rationale, proposed in this review's forward architecture, is essentially absent from the surveyed clinical literature and represents a concrete contribution.
\begin{table}[h]\centering\caption{Explainability method families.}\label{tab:xai}
\footnotesize\begin{tabular}{lll}\toprule
Family & Method & Caveat\\\midrule
Post-hoc attribution & SHAP, IG, Grad-CAM & faithfulness varies\\
Intrinsic & microstates, surrogate trees & lower ceiling\\
Attention maps & transformer weights & not causal explanation\\
Retrieval-grounded & RAG citations & absent in corpus\\\bottomrule
\end{tabular}\end{table}

\section{Benchmarks, Data Infrastructure, and Reproducibility}
Progress depends on shared infrastructure. CHB-MIT, TUH, Bonn, and Siena are the de-facto benchmarks; OmniEEG-Bench \citep{ep260600815} is an encouraging step toward standardised, leakage-aware evaluation. Yet reproducibility is weak: only about half the corpus names its dataset in the abstract, fewer release code, and per-subject results are rarely published. We argue the field needs a mandatory reporting standard---subject-wise splits, per-subject metrics, released code, and a model card---analogous to CONSORT in clinical trials.

\section{Historical Evolution, 2021--2026}
Table~\ref{tab:hist} traces the field's trajectory. The arc runs from CNN/RNN dominance (2021--22), through transformer adoption (2022--24), to graph, foundation, and state-space models (2024--26), with a parallel, persistent classical-baseline thread and a late-emerging consumer-grade and governance consciousness.
\begin{table}[h]\centering\caption{Field evolution by period.}\label{tab:hist}
\footnotesize\begin{tabular}{ll}\toprule
Period & Dominant development\\\midrule
2021--22 & CNN / RNN detectors, epoch-level reporting\\
2022--24 & transformers, multi-site assessment\\
2024--25 & graph NN, foundation models, domain adaptation\\
2025--26 & state-space (Mamba), consumer-grade, governance\\\bottomrule
\end{tabular}\end{table}

\section{Future Directions: A Detailed Agenda}
We specify seven concrete directions (Table~\ref{tab:future}), each falsifiable and tied to a gap. The throughline is that the next decade's advances will come less from new architectures than from \emph{honest evaluation}, \emph{cross-site robustness}, and \emph{governed deployment}.
\begin{table}[h]\centering\caption{Seven future directions.}\label{tab:future}
\footnotesize\begin{tabular}{cll}\toprule
\# & Direction & Gap\\\midrule
1 & Leakage-free community benchmark & G1\\
2 & Per-patient calibration / enrolment & G2\\
3 & Cross-dataset transfer protocols & G3\\
4 & Subject-wise consumer-grade benchmarks & G7\\
5 & Foundation models at clinical sensitivity & G12\\
6 & Fairness + calibration as requirements & G9\\
7 & Governed, auditable, HITL deployment & G10\\\bottomrule
\end{tabular}\end{table}

\section{Top-1\% Readiness Self-Assessment}
We assess this review against Q1 survey standards (Table~\ref{tab:q1}). It satisfies comprehensiveness (50 works, nine families), methodology (PRISMA), critical synthesis (the evaluation-honesty thesis), quantitative rigour (corpus statistics, $\chi^2$), forward contribution (this governed architecture architecture), and reproducibility (corpus + scripts released). Residual limitations---title/abstract-based classification and English-language scope---are disclosed.
\begin{table}[h]\centering\caption{Q1 survey-readiness checklist.}\label{tab:q1}
\footnotesize\begin{tabular}{ll}\toprule
Criterion & Status\\\midrule
$\geq$50 works, systematic search & yes (PRISMA)\\
Multi-axis taxonomy & yes (5 axes)\\
Per-paper comparison table & yes (47 rows, cited)\\
Critical thesis (not just catalogue) & yes (evaluation honesty)\\
Quantitative synthesis & yes ($\chi^2$, distributions)\\
Forward architecture & yes\\
Responsible/Explainable AI & yes (first-class)\\
Reproducibility & yes (corpus + scripts)\\
Limitations disclosed & yes\\\bottomrule
\end{tabular}\end{table}

\section{Comparison with Prior Surveys}
Prior epilepsy-EEG surveys catalogue architectures and datasets thoroughly but share three omissions this review addresses. First, they treat reported accuracy at face value rather than interrogating the epoch-vs-subject-wise distinction as a validity threat. Second, they stop at the model, omitting the deployment, governance, and regulatory layer that clinical translation requires. Third, they treat Responsible/Explainable/Interpretable AI as peripheral rather than as first-class, measurable reporting axes. By centring evaluation honesty, specifying a governed deployment architecture, and elevating responsible-AI metrics, this review is positioned as complementary to, and corrective of, the existing survey literature.

\section{Edge and Wearable Deployment}
Scalable screening outside specialist centres depends on edge and wearable deployment. Three hardware classes matter: clinical mobile EEG, consumer headsets (Emotiv-class), and minimal behind-the-ear or in-ear devices. Each trades electrode count and signal quality for accessibility. The modelling implication is severe: a 20-channel clinical detector cannot be transplanted to a 4-channel headset without retraining and, often, architectural change. Efficient architectures (pruned CNN, TCN, CNN-Mamba) and on-device quantisation are prerequisites. The honest expectation is lower sensitivity than clinical EEG, which makes the triage-with-escalation posture---not autonomous decision---the only defensible deployment for consumer hardware.

\section{Data Augmentation and Synthetic EEG}
Label scarcity and class imbalance motivate augmentation: time-warping, channel dropout, mixup, and frequency masking, plus generative synthesis (GANs, diffusion, VAEs) for minority seizure segments. Augmentation improves robustness but carries a subtle risk: synthetic segments that leak the statistics of the test subject reintroduce the very leakage subject-wise evaluation is meant to exclude. Augmentation policies must therefore be defined and applied within the training fold only. Synthetic EEG is promising for privacy-preserving data sharing but remains unvalidated at clinical sensitivity targets under subject-wise protocols.

\section{Active Learning and Annotation Efficiency}
Expert EEG annotation is the field's scarcest resource. Active learning---querying the expert only for the most informative segments---and weak/semi-supervised labelling reduce the annotation burden. These methods are particularly synergistic with the governance architecture: the same confidence and uncertainty estimates that drive human-in-the-loop escalation also identify the segments whose labels would most improve the model, closing a deploy-and-improve loop. Annotation efficiency is thus not merely a training-time convenience but a continuous-learning mechanism for deployed systems.

\section{Uncertainty Quantification in Depth}
A safe escalation policy is only as good as its uncertainty estimates. Four families apply: Bayesian/MC-dropout (epistemic uncertainty), deep ensembles (predictive variance), conformal prediction (distribution-free coverage guarantees), and evidential deep learning (single-pass uncertainty). Conformal prediction is especially attractive clinically because it provides finite-sample coverage guarantees without distributional assumptions---a property regulators value. The corpus almost never reports calibrated uncertainty, which is the single most important missing ingredient for the governed deployment this review advocates: without trustworthy uncertainty, confidence-tiered escalation is unsafe.

\section{Closed-Loop Neurostimulation and Reinforcement Learning}
Beyond passive detection, responsive neurostimulation (RNS) and deep-brain stimulation devices act on detected ictal activity in real time, closing the loop between sensing and therapy. This setting imposes the strictest constraints in the field: ultra-low latency, on-device inference, and an extremely low false-positive tolerance (every false trigger delivers unnecessary stimulation). Reinforcement-learning formulations---optimising stimulation policy against seizure-suppression reward---are emerging but face the same evaluation-honesty problem in amplified form: a policy validated on within-patient data may fail catastrophically on an unseen patient. Subject-wise, and ideally prospective, evaluation is non-negotiable for closed-loop systems, and the governance layer's kill-switch and audit row become safety-critical rather than merely advisable.

\section{Privacy-Preserving and Federated Learning}
EEG is sensitive health data; centralised aggregation across hospitals is often legally constrained. Federated learning \citep{ep250808159} trains a shared model without moving raw data, exchanging only gradients or weights; differential privacy and secure aggregation harden it against reconstruction attacks. Beyond compliance, federation directly serves generalisation: a model trained across many sites' distributions is, by construction, less prone to single-site over-fitting and the cross-dataset collapse documented above. The open challenges are communication efficiency, heterogeneity across site montages and labelling conventions, and the difficulty of subject-wise evaluation in a federated setting where no party sees all subjects. Federation is thus simultaneously a privacy mechanism and a generalisation mechanism.

\section{Paediatric versus Adult Epilepsy AI}
The field's dominant scalp benchmark, CHB-MIT, is paediatric; much clinical demand is adult. Seizure morphology, background rhythms, and artefact profiles differ substantially between the two populations, so a model validated on paediatric data does not transfer trivially to adults, and vice versa. This is a specific instance of the fairness and external-validity concerns raised above, and it has a concrete consequence: papers should state the age distribution of their data and avoid implying population-general performance from a single-cohort benchmark. Cross-population evaluation---train paediatric, test adult---is a revealing and rarely-run stress test that the field should adopt.

\section{Synthesis and Outlook}
Drawing the threads together: the epilepsy-EEG field has achieved remarkable within-recording accuracy and a rich architectural toolkit, but its central, under-acknowledged problem is that most reported performance does not survive honest, patient-independent evaluation. The corrective is not primarily architectural. It is methodological---subject-wise evaluation, cross-dataset testing, calibrated uncertainty, fairness reporting---and infrastructural---governed, explainable, auditable deployment with human oversight. This review has named that problem, quantified it, catalogued the gaps, surveyed every work, and specified a forward architecture that closes the deployment and integration gaps while demanding the methodological shift. The field's next decade will be defined less by which model wins a leaderboard than by whether the community adopts honest evaluation and deployable governance as default practice.

\section{The Leaderboard Critique}
Benchmark leaderboards have accelerated progress but also entrenched the evaluation-honesty problem: when a single epoch-level accuracy number ranks methods, the incentive is to optimise that number, not deployable performance. A leaderboard that ranks by epoch-level accuracy on CHB-MIT rewards exactly the leakage this review critiques. We argue community benchmarks should rank by subject-wise sensitivity at fixed specificity, report per-subject spread, and require released code---turning the leaderboard from a driver of the problem into a driver of its solution. OmniEEG-Bench's leakage-aware design \citep{ep260600815} is a step in this direction.

\section{Open-Source Tooling and the Reproducibility Stack}
A mature field needs a shared tooling stack. EEG preprocessing (MNE, EEGLAB), deep-learning frameworks (PyTorch, TensorFlow), and emerging EEG-specific model zoos lower the barrier to entry and improve reproducibility. The gap is downstream: there is no standard, shared implementation of leakage-free subject-wise evaluation, per-subject reporting, calibration measurement, or governance instrumentation. We argue the community would benefit from a reference evaluation harness that makes the honest protocol the path of least resistance---because methodology adoption follows tooling availability.

\section{Ethical and Societal Considerations}
Deploying seizure-detection AI carries ethical weight beyond accuracy. A missed seizure can cause harm; a false alarm erodes trust and may trigger unnecessary intervention. Consumer-grade screening democratises access but risks over-diagnosis and anxiety without clinical follow-up. Data sensitivity demands privacy-preserving training and storage. Equity demands fairness across age, sex, ethnicity, and site. And autonomy demands that AI inform, not replace, clinical judgement---hence the human-in-the-loop posture this review advocates throughout. These are not peripheral concerns; under a responsible-AI framing they are design requirements with concrete technical correlates.

\section{Glossary of Key Terms}
\begin{table}[h]\centering\caption{Glossary.}\label{tab:gloss}
\footnotesize\begin{tabular}{ll}\toprule
Term & Meaning\\\midrule
Epoch-level CV & segments split randomly (leakage-prone)\\
LOSO & leave-one-subject-out (patient-independent)\\
Sensitivity & true-positive rate (seizures caught)\\
Specificity & true-negative rate (normal recognised)\\
ECE / Brier & calibration error metrics\\
PSI & population stability index (drift)\\
RAG & retrieval-augmented generation\\
MCP & Model Context Protocol\\
AIOps & AI operations / monitoring\\
HITL & human-in-the-loop\\
this governed architecture & the governed deployment framework herein\\\bottomrule
\end{tabular}\end{table}
